\chapter{Battery operation}
\label{battery-operation}

At this stage you should have an LED blinking when the MCU is powered
from the ST-link.  The ST-Link provides 3.3\,V to power the MCU but
eventually you need to power your board from a battery using a
buck-converter to produce 5\,V that is reduced to 3.3\,V by a linear
regulator.


\section{Powering from power supply}

\begin{itemize}
\item Disconnect the ST-Link.

\item Twist the wires for the battery connector to reduce their
  inductance.

\item Configure a bench-top power supply to provide 7\,V with a
  current limit of 100\,mA.

\item Check the current limit by turning on the power supply with the
  output shorted.  The voltage should drop to 0\,V and the current
  should be 100\,mA.

\item Turn off the power supply and connect it to the battery
  connector wires.

\item Turn on the power supply.  If the current limits, turn off the
  power supply immediately and look for shorts.  This is when those
  0\,ohm isolating resistors come in handy.

\item If your LED does not flash, check the 3.3\,V voltage to the MCU.

\item Check the switching noise on the output of the buck-converter
  using an oscilloscope.  You will need to zoom in on the trace to
  50\,mV per division.  If you have more than 50\,mV of switching
  noise, your may have problems with your radio reception.  You may
  have a poor layout or need more capacitance.

\end{itemize}


\section{Powering from batteries}

About a week before the race you will be provided with batteries to
use instead of the power supply.
